% ===========================================================================
% Danh mục thuật ngữ và chữ viết tắt
% ===========================================================================
\phantomsection
\addcontentsline{toc}{chapter}{DANH MỤC THUẬT NGỮ VÀ CHỮ VIẾT TẮT}
\chapter*{DANH MỤC THUẬT NGỮ VÀ CHỮ VIẾT TẮT}
\markboth{DANH MỤC THUẬT NGỮ VÀ CHỮ VIẾT TẮT}{}

\begingroup
\small
\setlength{\extrarowheight}{1.5pt}
\begin{longtable}{@{}L{2.2cm}L{5.0cm}L{6.6cm}@{}}
\toprule
\textbf{Viết tắt} & \textbf{Thuật ngữ tiếng Anh} & \textbf{Giải nghĩa tiếng Việt} \\
\midrule
\endfirsthead
\multicolumn{3}{@{}l}{\small\textit{(tiếp theo trang trước)}}\\
\toprule
\textbf{Viết tắt} & \textbf{Thuật ngữ tiếng Anh} & \textbf{Giải nghĩa tiếng Việt} \\
\midrule
\endhead
\bottomrule
\multicolumn{3}{r@{}}{\small\textit{(xem tiếp trang sau)}}\\
\endfoot
\bottomrule
\endlastfoot

API   & Application Programming Interface & Giao diện lập trình ứng dụng \\
AUC   & Area Under the ROC Curve & Diện tích dưới đường cong ROC \\
CART  & Classification and Regression Tree & Cây phân loại và hồi quy \\
CIC   & Credit Information Center & Trung tâm Thông tin tín dụng Quốc gia Việt Nam \\
DPD   & Days Past Due & Số ngày quá hạn \\
DTI   & Debt-to-Income ratio & Tỷ lệ nghĩa vụ trả nợ trên thu nhập \\
EAD   & Exposure at Default & Dư nợ tại thời điểm vỡ nợ \\
ECE   & Expected Calibration Error & Sai số hiệu chuẩn kỳ vọng \\
EL    & Expected Loss & Tổn thất kỳ vọng \\
FN    & False Negative & Âm tính giả (bỏ sót hồ sơ xấu) \\
FP    & False Positive & Dương tính giả (từ chối nhầm hồ sơ tốt) \\
GBM   & Gradient Boosting Machine & Máy tăng cường độ dốc \\
IV    & Information Value & Giá trị thông tin của biến \\
JSON  & JavaScript Object Notation & Định dạng trao đổi dữ liệu JSON \\
KS    & Kolmogorov--Smirnov statistic & Thống kê Kolmogorov--Smirnov \\
LGD   & Loss Given Default & Tỷ lệ tổn thất khi vỡ nợ \\
LR    & Logistic Regression & Hồi quy logistic \\
LTV   & Loan-to-Value ratio & Tỷ lệ vay trên giá trị tài sản bảo đảm \\
MLP   & Multi-Layer Perceptron & Mạng nơ-ron truyền thẳng nhiều lớp \\
NHNN  & --- & Ngân hàng Nhà nước Việt Nam \\
NPL   & Non-Performing Loan & Nợ xấu (nợ nhóm 3--5) \\
PAVA  & Pool Adjacent Violators Algorithm & Thuật toán gộp vi phạm liền kề \\
PD    & Probability of Default & Xác suất vỡ nợ \\
PR-AUC & Area Under Precision--Recall Curve & Diện tích dưới đường Precision--Recall \\
PSI   & Population Stability Index & Chỉ số ổn định phân phối \\
PTI   & Payment-to-Income ratio & Tỷ lệ khoản trả góp trên thu nhập \\
RBF   & Radial Basis Function & Hàm cơ sở bán kính (hạt nhân Gauss) \\
REST  & Representational State Transfer & Kiến trúc dịch vụ web REST \\
RF    & Random Forest & Rừng ngẫu nhiên \\
ROC   & Receiver Operating Characteristic & Đường đặc trưng hoạt động của bộ thu nhận \\
SHAP  & SHapley Additive exPlanations & Phương pháp giải thích mô hình dựa trên giá trị Shapley \\
SQL   & Structured Query Language & Ngôn ngữ truy vấn có cấu trúc \\
SVM   & Support Vector Machine & Máy vectơ hỗ trợ \\
TCTD  & --- & Tổ chức tín dụng \\
WOE   & Weight of Evidence & Trọng số bằng chứng \\
XGBoost & eXtreme Gradient Boosting & Thư viện tăng cường độ dốc cực hạn \\
\end{longtable}
\endgroup

\vspace{0.5cm}
\noindent\textbf{Một số thuật ngữ được sử dụng thống nhất trong đề án:}

\begingroup
\small
\begin{longtable}{@{}L{4.4cm}L{9.6cm}@{}}
\toprule
\textbf{Thuật ngữ} & \textbf{Cách hiểu trong đề án} \\
\midrule
\endfirsthead
\toprule
\textbf{Thuật ngữ} & \textbf{Cách hiểu trong đề án} \\
\midrule
\endhead
\bottomrule
\endlastfoot

Kiểm duyệt tín dụng &
Quá trình đánh giá, sàng lọc hồ sơ vay để hình thành đề xuất có căn cứ. Phân biệt
với \textit{phê duyệt} -- hành vi ra quyết định của cấp có thẩm quyền. \\

Mô hình hỗ trợ kiểm duyệt &
Mô hình cung cấp đánh giá rủi ro định lượng kèm lý giải làm đầu vào cho quyết
định của con người, không thay thế thẩm quyền phê duyệt. \\

Tính đúng đắn theo thời điểm \newline (\eng{point-in-time correctness}) &
Nguyên tắc bảo đảm mọi đặc trưng của một hồ sơ chỉ được tính từ dữ liệu đã tồn
tại trước hoặc tại thời điểm nộp hồ sơ. \\

Rò rỉ nhãn \newline (\eng{label leakage}) &
Hiện tượng đặc trưng chứa thông tin phát sinh sau thời điểm nộp hồ sơ, khiến mô
hình cho kết quả cao giả tạo khi đánh giá và sụp đổ khi triển khai. \\

Cửa sổ quan sát \newline (\eng{performance window}) &
Khoảng thời gian kể từ khi cấp tín dụng dùng để xác định hồ sơ có phát sinh nợ
xấu hay không. Đề án sử dụng cửa sổ 12 tháng. \\

Kiểm định ngoài thời gian \newline (\eng{out-of-time validation}) &
Cách đánh giá mô hình trên dữ liệu thuộc giai đoạn sau giai đoạn huấn luyện, phản
ánh đúng điều kiện vận hành thực tế. \\

Trần Bayes &
Giới hạn lý thuyết về chất lượng dự báo mà không mô hình nào vượt qua được, tính
từ xác suất vỡ nợ thật của quá trình sinh dữ liệu. \\

Hiệu chuẩn xác suất \newline (\eng{calibration}) &
Quá trình điều chỉnh để giá trị xác suất mô hình đưa ra trùng khớp với tần suất
xảy ra thực tế -- điều kiện cần để dùng PD trong phân vùng quyết định. \\

Khách hàng thin-file &
Khách hàng có rất ít hoặc không có lịch sử tín dụng tại CIC, khiến việc đánh giá
phải dựa chủ yếu vào dữ liệu hành vi giao dịch. \\

Cắt ngọn \newline (\eng{winsorize}) &
Kỹ thuật giới hạn giá trị cực trị về mức phân vị biên thay vì loại bỏ quan sát,
nhằm giữ nguyên cỡ mẫu. \\

Trôi dạt dữ liệu \newline (\eng{data drift}) &
Hiện tượng phân phối dữ liệu hoặc quan hệ giữa đặc trưng và nhãn thay đổi theo
thời gian, làm mô hình suy giảm chất lượng. \\
\end{longtable}
\endgroup
